\chapter{Giơ tay chọn mức tỉnh}
\begin{tinhhuongbox}{Tình huống | Situation}
\songngu{Đầu giờ sáng hoặc đầu giờ chiều, lớp thường chưa vào cùng một trạng thái.}{At the start of the morning or the afternoon, students are rarely in the same state.}
\songngu{Thầy cô cần đo nhanh bầu không khí trước khi quyết định nhịp dạy.}{You need a quick read of the room before deciding the teaching pace.}
\end{tinhhuongbox}
\begin{noitrenlop}
\songngu{Ai đang tỉnh thì giơ một tay, ai cần cứu nhẹ thì giơ hai tay, đừng giơ ba tay kẻo thành phim kinh dị.}{If you are awake, raise one hand. If you need a little rescue, raise two. Do not raise three, or this will turn into a horror movie.}
\end{noitrenlop}
\begin{vihieuqua}
\songngu{Cả lớp cùng phản hồi trong vài giây và tiếng cười làm không khí mở ra rất nhanh.}{The whole class responds in a few seconds, and the laugh opens the room very quickly.}
\end{vihieuqua}
\begin{englishpocket}
\songngu{Cụm nên nhớ: \textit{I need a little rescue.}}{Phrase to keep: \textit{I need a little rescue.}}
\end{englishpocket}
\chapter{Ba từ để vào bài}
\begin{tinhhuongbox}{Tình huống | Situation}
\songngu{Muốn khởi động nhanh mà không cần đạo cụ, tốt nhất là dùng giới hạn nhỏ.}{If you want a quick warm-up without any materials, the best move is to use a small limit.}
\songngu{Ba từ đủ ít để ai cũng dám nói, nhưng đủ nhiều để suy nghĩ bật lên.}{Three words are few enough for anyone to say, but many enough to trigger thinking.}
\end{tinhhuongbox}
\begin{noitrenlop}
\songngu{Mỗi bàn ba từ thôi để mở bài, ai nói dài quá tôi thu học phí phát biểu thêm.}{Each desk gets only three words to open the lesson. If anyone speaks too long, I will charge extra speaking tuition.}
\end{noitrenlop}
\begin{vihieuqua}
\songngu{Giới hạn ngắn làm lớp gọn lại mà không bị ép im.}{A short limit makes the class more concise without forcing silence.}
\end{vihieuqua}
\begin{englishpocket}
\songngu{Cụm nên nhớ: \textit{Only three words.}}{Phrase to keep: \textit{Only three words.}}
\end{englishpocket}
\chapter{Bắt đầu bằng một ví dụ gần lớp}
\begin{tinhhuongbox}{Tình huống | Situation}
\songngu{Bài học xa quá thì đầu tiết sẽ khó kéo.}{If the lesson feels too distant, the start will be hard to pull off.}
\songngu{Ví dụ gần lớp là cây cầu nhanh nhất từ đời sống sang nội dung.}{A nearby example is the fastest bridge from life to content.}
\end{tinhhuongbox}
\begin{noitrenlop}
\songngu{Trước khi học cái lớn, ai cho tôi một ví dụ ngay trong lớp này để đỡ phải bay quá cao?}{Before we study the big thing, who can give me one example from this very room so we do not have to fly too high too soon?}
\end{noitrenlop}
\begin{vihieuqua}
\songngu{Học sinh thấy bài học gần mình hơn nên dễ nhập hơn.}{Students feel the lesson is closer to them, so they enter it more easily.}
\end{vihieuqua}
\begin{englishpocket}
\songngu{Cụm nên nhớ: \textit{From this very room.}}{Phrase to keep: \textit{From this very room.}}
\end{englishpocket}
\chapter{Một cặp đôi mở tiết}
\begin{tinhhuongbox}{Tình huống | Situation}
\songngu{Có lớp ngại phát biểu cá nhân nhưng lại nói tốt hơn khi quay sang bạn bên cạnh.}{Some classes dislike speaking alone but do much better when they turn to the person next to them.}
\songngu{Khởi động theo cặp làm mức an toàn tăng lên ngay.}{A pair-based warm-up immediately raises the sense of safety.}
\end{tinhhuongbox}
\begin{noitrenlop}
\songngu{Quay sang bạn bên cạnh, mỗi người nói một câu thôi, rồi tôi gọi ngẫu nhiên để cứu sự run của cả hai.}{Turn to the person next to you, each of you says just one sentence, and then I will call on a random pair to save both of you from panic.}
\end{noitrenlop}
\begin{vihieuqua}
\songngu{Nói trước với một người giúp học sinh dễ nói trước cả lớp hơn.}{Speaking to one person first helps students speak to the whole class more easily.}
\end{vihieuqua}
\begin{englishpocket}
\songngu{Cụm nên nhớ: \textit{Turn to the person next to you.}}{Phrase to keep: \textit{Turn to the person next to you.}}
\end{englishpocket}
\chapter{Một câu hẹn cho cuối tiết}
\begin{tinhhuongbox}{Tình huống | Situation}
\songngu{Khởi động đẹp nhất là khởi động mở đường cho cả tiết, không chỉ làm lớp cười một chút rồi thôi.}{The best warm-up opens the path for the whole lesson, not just a quick laugh.}
\songngu{Một câu hẹn đầu giờ sẽ giữ sợi dây tới cuối giờ.}{One promise at the start can hold a thread until the end.}
\end{tinhhuongbox}
\begin{noitrenlop}
\songngu{Đầu giờ tôi hỏi một câu thôi, cuối giờ quay lại xem lớp trả lời hay hơn được chưa nhé.}{I will ask only one question now, and at the end we will come back and see whether the class can answer it better.}
\end{noitrenlop}
\begin{vihieuqua}
\songngu{Lớp có lý do để đi hết tiết thay vì chỉ đi hết từng đoạn nhỏ.}{The class gets a reason to travel through the whole lesson instead of only surviving one small segment at a time.}
\end{vihieuqua}
\begin{englishpocket}
\songngu{Cụm nên nhớ: \textit{We will come back to it.}}{Phrase to keep: \textit{We will come back to it.}}
\end{englishpocket}